The metrics induced by AutoLibra (Table \ref{tab:metrics}) capture the behavior of the agent effectively, with coverage increasing from 65\% at Iteration 0 to 92\% at \emph{Iteration 3} and a low mean redundancy of 56\%. Notably, earlier metrics were found to describe higher-level behaviors than later metrics, reflective of the more complex behaviors that the agent demonstrated in later iterations, as well as the compositionality of induced metrics. Code changes were selected to specifically target given metrics, and as seen in Figure \ref{fig:autolibra-training}, the agent's performance on the targeted metric improved significantly in the iteration following the code change. This demonstrates the utility of AutoLibra for fine-grained agent improvement, as well as the human interpretability of the induced metrics.

The section below discusses iteration-by-iteration induced metrics, code changes directed by the metrics, and the results of these code changes on agent environment and trajectory performance.

\textbf{Iteration 0}
The agent behavior is stochastic and inconsistent, with no clear strategy or goal. Any progress in the environment is the consequence of a random walk. 
\texttt{Win Condition Recognition} \includegraphics[scale=0.07]{figs/emojis/emoji_1.png} was identified as a prerequisite to all other metrics induced in this iteration, as progress in the environment is impossible without recognizing the win condition, and was therefore targeted for improvement. Improvements took the form of few-shot prompting, by supplying an example of observations and the corresponding win condition and subtasks, and guidance on subtask-level planning, by mapping specific environmental observations to actions that would progress the agent towards the win condition.
\textbf{Iteration 1}
An increase in \includegraphics[scale=0.07]{figs/emojis/emoji_1.png} performance of 22\% was observed between \emph{Iteration 0} and \emph{Iteration 1}, indicating that the code changes were effective in improving the agent's ability to recognize the win condition. A slight but statistically insignificant increase was observed for the other metrics induced in \emph{Iteration 0}, but the agent's overall baba-is-ai environment score increased by 10\%, indicating that \includegraphics[scale=0.07]{figs/emojis/emoji_1.png} was correctly identified by AutoLibra as a key bottleneck to the agent's performance in the environment.

The agent now targets specific blocks and objects instead of moving randomly, but gets stuck in loops and on immovable objects, and is unable to consistently complete multi-step tasks.

Based on the metrics induced in \emph{Iteration 1}, several changes to the agent code were implemented:
\begin{itemize}
    \item Augmentation of existing subtask-level planning guidance with low-level position instructions, targeting improvement of \texttt{Rule Modification for Obstacle Management} \includegraphics[scale=0.07]{figs/emojis/emoji_2.png} and \texttt{Direct Navigation Efficiency} \includegraphics[scale=0.07]{figs/emojis/emoji_3.png}
    \item Meta-prompting \cite{Suzgun_Kalai_2024} by providing subtask identification heuristic, targeting improvement of \texttt{Selective Interaction with Relevant Objects} \includegraphics[scale=0.07]{figs/emojis/emoji_6.png} and \texttt{Context-Sensitive Decision Making} \includegraphics[scale=0.07]{figs/emojis/emoji_4.png}
    \item Augmentation of single-shot example with movement instructions, targeting improvement of \texttt{Rule Manipulation and Execution} \includegraphics[scale=0.07]{figs/emojis/emoji_7.png}
\end{itemize}


\textbf{Iteration 2}
Substantial improvements in \includegraphics[scale=0.07]{figs/emojis/emoji_1.png} were observed due to more detailed single-shot examples and reasoning templates, and the corresponding increase in other metrics supports idea that \includegraphics[scale=0.07]{figs/emojis/emoji_1.png} is a key bottleneck to the agent's performance. We further observed a significant increase in \includegraphics[scale=0.07]{figs/emojis/emoji_2.png} and \includegraphics[scale=0.07]{figs/emojis/emoji_7.png}, indicating that the changes targeting these metrics were effective in improving the agent's reasoning performance and consequent avoidance of irrelevant objects. \texttt{Subtask Identification} \includegraphics[scale=0.07]{figs/emojis/emoji_5.png} saw no significant improvement, as the agent still misunderstands map directions and confuses rule blocks with objects. 

Qualitatively, the agent now recognizes when wall rules need to be broken and breaks them, but still cannot assemble rules to win, as it gets stuck on immovable blocks and walls and does not understand where rule blocks need to be placed to form a win condition.

Based on the metrics induced in \emph{Iteration 2}, several changes to the agent code were implemented:
\begin{itemize}
    \item Formatting observations as absolute position (as opposed to relative steps from the agent), and listing of immovable/movable blocks to improve \texttt{Direct Navigation Efficiency} \includegraphics[scale=0.07]{figs/emojis/emoji_3.png}
    \item Chain-of-thought reasoning to assemble subtask completion agenda based on iteration-to-iteration observations, targeting improvement of \texttt{Subtask Coordination and Overall Task Planning} \includegraphics[scale=0.07]{figs/emojis/emoji_8.png}
    \item Augmentation of single-shot example to few-shot example with reasoning templates, targeting improvement of \texttt{Rule Manipulation and Execution} \includegraphics[scale=0.07]{figs/emojis/emoji_7.png}
    \item Explicit instructions on navigation to avoid block collisions and assembly of rules, including movement templates, targeting improvement of \texttt{Rule Manipulation and Execution} \includegraphics[scale=0.07]{figs/emojis/emoji_7.png} and \texttt{Win Rule Construction} \includegraphics[scale=0.07]{figs/emojis/emoji_5.png}
\end{itemize}

\textbf{Iteration 3}
We observe a large increase in \includegraphics[scale=0.07]{figs/emojis/emoji_4.png} and \includegraphics[scale=0.07]{figs/emojis/emoji_6.png}, indicating that the changes targeting these metrics were effective in improving the agent's reasoning performance and consequent avoidance of irrelevant objects. \includegraphics[scale=0.07]{figs/emojis/emoji_3.png} saw further improvement, indicating that the use of absolute position and immovable/movable block listing was effective in improving the agent's navigation efficiency. \includegraphics[scale=0.07]{figs/emojis/emoji_7.png} saw a slight increase, indicating that the changes targeting this metric were effective in improving the agent's ability to manipulate rules to achieve the win condition. \includegraphics[scale=0.07]{figs/emojis/emoji_5.png} saw its first improvement, indicating that more complex metrics are effectively targeted when the base performance and simpler metrics are improved.

Qualitatively, nearly every agent always understands when wall rules can be and cannot be broken, as well as when it's not necessary to break the wall rule, covers nearly all tasks that involve going to win. The agent understands how to assemble win rules but still struggles with changing block direction and understanding that blocks get stuck against walls.